\section{Exemplos resolvidos}

\begin{frame}{Exemplo 2 -- alimentador trifásico}
\small
\begin{columns}[T,onlytextwidth]
\column{0.50\textwidth}
\begin{caixaEx}
$P=\SI{42}{kW}$, $V=\SI{380}{V}$, $fp=0{,}88$, $L=\SI{60}{m}$.
\[
I_b=\frac{42000}{\sqrt3\cdot380\cdot0{,}88}\approx\SI{72.5}{A}.
\]
\end{caixaEx}
\column{0.48\textwidth}
\begin{caixaProjeto}[title=Sequência]
\begin{enumerate}
\item selecionar seção por ampacidade corrigida;
\item selecionar proteção;
\item verificar $\Delta V$;
\item verificar curto máximo/mínimo;
\item verificar PE e neutro;
\item revisar seletividade a montante/jusante.
\end{enumerate}
\end{caixaProjeto}
\end{columns}
\end{frame}
\begin{frame}{Exemplo 3 -- banco de capacitores}
\small
\begin{columns}[T,onlytextwidth]
\column{0.50\textwidth}
\begin{caixaEx}
$P=\SI{60}{kW}$, corrigir $fp$ de $0{,}80$ para $0{,}95$.
\begin{align*}
\varphi_1&=36{,}87^\circ,\quad \varphi_2=18{,}19^\circ\\
Q_c&=60(\tg\varphi_1-\tg\varphi_2)\\
&\approx\SI{25.3}{kvar}.
\end{align*}
\end{caixaEx}
\column{0.48\textwidth}
\begin{caixaAt}
Antes de instalar $\SI{25}{kvar}$ ``puro'', avaliar:
\begin{itemize}
\item perfil de carga;
\item risco de sobrecompensação;
\item harmônicas;
\item ressonância;
\item número de estágios;
\item local de compensação.
\end{itemize}
\end{caixaAt}
\end{columns}
\end{frame}
\begin{frame}{Exemplo 4 -- motor de 22 kW em rede fraca}
\small
\begin{columns}[T,onlytextwidth]
\column{0.50\textwidth}
\begin{caixaEx}
$P=\SI{22}{kW}$, $V=\SI{380}{V}$, $fp=0{,}86$, $\eta=0{,}93$.
\[
I_N\approx\frac{22000}{\sqrt3\cdot380\cdot0{,}86\cdot0{,}93}
\approx\SI{41.8}{A}.
\]
\end{caixaEx}
\column{0.48\textwidth}
\begin{caixaProjeto}[title=Decisão]
Se a rede não tolera inrush elevado:
\begin{itemize}
\item calcular queda durante partida;
\item verificar torque requerido;
\item comparar soft-starter e VFD;
\item coordenar proteção;
\item verificar harmônicas/EMC.
\end{itemize}
\end{caixaProjeto}
\end{columns}
\end{frame}
\begin{frame}{Exemplo 5 -- curto no secundário do transformador}
\small
\begin{columns}[T,onlytextwidth]
\column{0.51\textwidth}
\begin{caixaEx}
Trafo $\SI{500}{kVA}$, $\SI{380}{V}$, $Z=5{,}5\%$.
\[
I_N=\frac{500000}{\sqrt3\cdot380}\approx\SI{760}{A}.
\]
\[
I_{cc}\approx\frac{760}{0{,}055}\approx\SI{13.8}{kA}.
\]
\end{caixaEx}
\column{0.47\textwidth}
\begin{caixaObs}
No QGBT real, incluir impedância da rede de MT, trafo, conexões e barramentos. Nos QDs a jusante, incluir cabos/alimentadores, reduzindo $I_{cc}$.
\end{caixaObs}
\end{columns}
\end{frame}
\begin{frame}{Exemplo 6 -- balanceamento de fases}
\small
\begin{columns}[T,onlytextwidth]
\column{0.50\textwidth}
\begin{caixaEx}
Correntes medidas:
\[
I_A=\SI{52}{A},\quad I_B=\SI{38}{A},\quad I_C=\SI{44}{A}.
\]
\[
I_{med}=\SI{44.7}{A}.
\]
A fase A está $\approx\SI{7.3}{A}$ acima da média.
\end{caixaEx}
\column{0.48\textwidth}
\begin{caixaProjeto}[title=Ação]
Mover circuitos monofásicos de A para B/C, mas respeitando:
\begin{itemize}
\item simultaneidade real;
\item cargas críticas;
\item sequência de fases;
\item DRs e neutros independentes;
\item harmônicas no neutro.
\end{itemize}
\end{caixaProjeto}
\end{columns}
\end{frame}

